\section{导言：为什么最后一讲不再讲新算法}

这一讲不是在课程末尾再补一套新模型，而是把整门 CS231N 的技术线索重新放回到一个更大的问题里：\textbf{视觉 AI 到底应该学会什么，以及这些能力最终服务于谁。} 讲者一开始就明确说明，这是一场``perspective talk''，目标不是引入新 loss、新网络结构，而是从研究史、社会影响和 human perspective 三个维度，把已经学过的视觉技术重新组织起来。\footnote{根据 lecture18 字幕 00:00:05--00:01:15，讲者明确说明本讲不再教授新的算法材料，而是从 long-term research evolution 与 human perspective 回看 AI。}

\begin{importantbox}{本讲不是算法加餐，而是课程总收束}
\begin{enumerate}
    \item 它回顾视觉智能的发展主线：从人类视觉能力启发，到对象识别、关系理解、视频理解。
    \item 它强调 AI 的第二层价值：不仅复制人类能力，还要帮助人类看见看不到、看不全、看不稳的东西。
    \item 它把最终问题拉回到价值对齐：AI 不只是``会不会做''，还包括``该不该做''、``怎样做才尊重人''。
\end{enumerate}
\end{importantbox}

从结构上看，整场讲座可以压缩成三句话：

\begin{itemize}
    \item \textbf{See what humans see}：先理解人类为什么能把视觉变成一种高效的智能能力。
    \item \textbf{See what humans don't see}：再讨论 AI 如何超越或补足人类感知的盲点。
    \item \textbf{See what humans want to see}：最后约束这些能力，让它们服务于人的真实偏好与价值。
\end{itemize}

\begin{figure}[H]
\centering
\begin{tikzpicture}[node distance=1.35cm, every node/.style={font=\small}]
\node[draw, rounded corners, fill=blue!8, align=center, minimum width=4.8cm, minimum height=0.95cm] (a) {\textbf{See what humans see}\\复制并理解人类已有视觉能力};
\node[draw, rounded corners, fill=yellow!16, align=center, below=of a, minimum width=4.8cm, minimum height=0.95cm] (b) {\textbf{See what humans don't see}\\超越盲点、稳定性与尺度限制};
\node[draw, rounded corners, fill=green!14, align=center, below=of b, minimum width=4.8cm, minimum height=0.95cm] (c) {\textbf{See what humans want to see}\\把能力约束到人的目标与边界};
\draw[-{Latex[length=3mm]}, thick] (a) -- (b);
\draw[-{Latex[length=3mm]}, thick] (b) -- (c);
\end{tikzpicture}
\caption{Lecture 18 的三段式逻辑：从能力、到补足、再到价值对齐}
\end{figure}

\begin{knowledgebox}{读这讲的正确方式}
这不是一篇``技术细节最密''的 lecture，而是一篇``概念跨度最大''的 lecture。它的价值不在于证明某个模型更强，而在于给出一个判断框架：当我们构建视觉系统时，应该同时问\textbf{能力边界、应用场景、人的偏好、风险与责任}。
\end{knowledgebox}

